\documentclass[12pt,a4paper]{ctexart}

\usepackage[margin=2.5cm]{geometry}
\usepackage{hyperref}

\hypersetup{
  hidelinks
}

\setlength{\parskip}{0.45em}
\setlength{\parindent}{2em}

\begin{document}

\begin{center}
{\LARGE\heiti 测试报告}
\end{center}

\tableofcontents
\newpage

\section{简介}

\subsection{编写目的}

本报告用于记录 \texttt{RAG-Knowledge-Assistant} 项目的测试设计、测试执行情况、测试结果分析及后续扩展适配方案，作为课程作业中测试工作的最终总结材料。报告重点说明黑盒测试与白盒测试的测试范围、测试方法、执行结果、缺陷情况和成员贡献，便于对系统当前质量和测试充分性进行综合评价。

\subsection{项目背景}

\texttt{RAG-Knowledge-Assistant} 是面向课程学习场景的知识辅助系统，主要包含用户与认证、课程管理、课程资料管理、练习管理和课程社区管理等功能模块。系统采用前后端分离架构，前端负责页面交互与业务展示，后端提供课程、资料、练习和社区相关接口，并结合文档存储与关系型数据存储完成核心业务支撑。

\subsection{测试范围}

本次测试围绕课程作业要求和系统已实现功能展开，主要包括黑盒功能测试与白盒代码测试两部分。黑盒测试从用户使用视角验证注册登录、课程创建与加入、资料上传与预览、练习题管理、课程社区互动以及核心业务流程；白盒测试从代码结构视角验证用户与认证、课程管理、资料管理、练习管理和社区管理相关后端模块的语句覆盖、分支覆盖、基本路径和关键数据流。性能测试、压力测试、故障注入测试和专项安全测试未作为本次课程作业的主要测试范围。

\section{黑盒测试报告}

\subsection{测试概述}

本次黑盒测试面向 \texttt{RAG-Knowledge-Assistant} 项目展开，围绕用户管理、课程管理、课程资料管理、练习题管理、课程社区管理以及核心业务流程，对系统主要业务功能进行验证。

测试用例设计主要采用等价类划分、边界值分析、决策表法和场景法。测试执行采用“自动化接口测试为主、少量人工页面测试为辅”的方式，其中自动化测试承担主体覆盖任务，人工测试用于补充前端表单交互、PDF 预览、下载行为和页面级端到端流程验证，并结合实际运行结果与人工验证记录进行综合判断。

\subsection{测试环境与工具}

\subsubsection{测试环境}

\begin{enumerate}
  \item 操作系统：Windows 11 开发环境。
  \item 前端：React 19。
  \item 后端：FastAPI。
  \item 浏览器：Chrome 与 Edge。
  \item 前端访问地址：\url{http://localhost:3000}。
  \item 后端接口地址：\url{http://127.0.0.1:8000}。
  \item 自动化测试环境：SQLite 测试库用于承载关系型业务数据；资料模块使用文档存储替身模拟文件行为。
\end{enumerate}

\subsubsection{测试工具}

\begin{enumerate}
  \item \texttt{Codex}：用于辅助生成测试计划、测试用例、自动化测试代码和测试报告整理内容。
  \item \texttt{pytest + FastAPI TestClient}：用于执行自动化黑盒接口测试。
  \item 浏览器人工操作：用于执行少量页面黑盒验证。
\end{enumerate}

\subsection{测试执行结果}

\subsubsection{自动化测试结果}

自动化黑盒测试共执行 36 项，全部通过。

执行命令：

\begin{verbatim}
D:\miniconda3\envs\rag\python.exe -m pytest -q
\end{verbatim}

执行结果：

\begin{verbatim}
36 passed, 9 warnings
\end{verbatim}

自动化测试已覆盖用户管理、课程管理、资料管理、练习题管理、课程社区管理以及核心流程的主要黑盒接口验证，具体包括注册、登录、个人信息查看与修改、课程创建、加入与退出、资料上传、查看与删除、练习题查看与作答、社区发帖、评论、回复等主要功能点。

\subsubsection{人工测试结果}

人工黑盒测试共执行 10 项，全部通过。

人工测试主要补充验证了以下内容：

\begin{enumerate}
  \item 注册页确认密码校验与角色选择校验
  \item 登录后首页跳转、身份状态展示与刷新后的登录状态保持
  \item PDF 资料预览效果与页面展示情况
  \item 资料下载时的浏览器下载行为
  \item 练习题页面反馈展示与“重新作答”后的状态恢复
  \item 匿名发帖与匿名评论的展示效果
  \item 回复评论后的页面定位与回显
  \item 页面级核心流程端到端验证
\end{enumerate}

\subsubsection{综合统计}

\begin{enumerate}
  \item 自动化测试：36 项，通过 36 项。
  \item 人工测试：10 项，通过 10 项。
  \item 黑盒测试总执行数：46 项。
  \item 黑盒测试总通过数：46 项。
\end{enumerate}

自动化测试与人工测试结果一致，计划范围内的主要黑盒功能均完成验证。

\subsection{缺陷与问题分析}

本次黑盒测试未发现阻塞主要业务流程的严重缺陷，系统主要功能能够按需求正常运行，未出现导致系统无法使用、核心流程中断或角色权限错误的问题。

测试过程中记录到的轻微问题或风险如下：

\begin{enumerate}
  \item 回复评论定位功能在页面滚动距离较长时，定位体验可能不够稳定。该问题未导致功能失效，因此未判定为失败项，但建议在后续版本中优化页面滚动与定位交互体验。
  \item 自动化测试执行中出现若干技术栈弃用警告。相关警告主要来自 Pydantic、SQLAlchemy 和 FastAPI 的版本兼容提示，不影响本次黑盒测试通过结果，但说明项目后续存在升级维护需求。
\end{enumerate}

\subsection{测试结论}

本次黑盒测试结果表明，\texttt{RAG-Knowledge-Assistant} 项目在当前测试范围内整体表现良好，主要业务功能满足需求说明文档要求，教师与学生角色权限控制正确，课程平台核心业务流程能够顺利完成。

综合自动化测试与人工测试结果，可以认为该系统在黑盒测试范围内达到了课程作业要求，具备较好的功能完整性和可用性。

\section{白盒测试报告}

\subsection{测试概述}

本次白盒测试重点覆盖用户与认证、课程管理、课程资料管理、练习管理和课程社区管理五类业务。测试目标是基于代码结构验证关键逻辑的正确性与稳健性，并满足课程作业对白盒测试方法与覆盖率指标的要求。

白盒测试的执行策略：优先通过后端接口触发真实业务调用链，贯通 \texttt{controller}、\texttt{service}、\texttt{repository}、\texttt{dao} 等后端层次，仅对外部系统依赖（如文档存储客户端）使用替身隔离，以保证测试结果既具可重复性又具工程真实性。

本阶段白盒测试由成员 B 与成员 C 分工完成。成员 B 负责用户与认证、课程管理、课程资料管理相关代码；成员 C 负责练习管理与课程社区管理相关代码。两部分均采用代码覆盖分析、基本路径覆盖和数据流分析三类方法，测试组织方式和结果统计口径保持一致。

\subsection{测试环境与工具}

\subsubsection{测试环境}

\begin{enumerate}
  \item 运行环境：Python 3.11 / Python 3.12。
  \item 后端框架：FastAPI。
  \item 关系型数据测试环境：SQLite 内存数据库（测试会话）。
  \item 文档存储测试环境：Mongo/GridFS 外部依赖替身。
\end{enumerate}

\subsubsection{测试工具}

\begin{enumerate}
  \item \texttt{Codex}：辅助设计测试用例、分析未覆盖路径和整理报告。
  \item \texttt{pytest}：执行自动化白盒测试。
  \item \texttt{FastAPI TestClient}：通过接口请求驱动后端完整调用链。
  \item \texttt{pytest-cov}：统计覆盖率并输出未命中行。
\end{enumerate}

\subsection{测试对象说明}

白盒测试覆盖如下代码模块：

\begin{enumerate}
  \item 用户与认证调用链：\texttt{LoginController \(\rightarrow\) LoginService \(\rightarrow\) UserRepository + UserAuthRepository}；\texttt{UserController \(\rightarrow\) UserService \(\rightarrow\) UserRepository + UserAuthRepository}。
  \item 课程管理调用链：\texttt{CourseController \(\rightarrow\) CourseService \(\rightarrow\) CourseRepository + UserRepository}。
  \item 资料管理调用链：\texttt{rag\_controller \(\rightarrow\) DocumentDAO \(\rightarrow\) document\_repository}。
  \item 练习管理调用链：\texttt{ExerciseController}、\texttt{ExerciseService} 与 \texttt{ExerciseRepository}。
  \item 课程社区调用链：\texttt{PostController}、\texttt{PostService}、\texttt{PostRepository} 与 \texttt{CommentRepository}。
\end{enumerate}

\subsection{测试方法说明}

本阶段白盒测试主要采用代码覆盖分析、基本路径覆盖和数据流分析三种方法。结合本项目后端调用链特点，各方法表现出如下特点：

\subsubsection{代码覆盖分析特点}

\begin{enumerate}
  \item 优点：可量化评估测试充分性，能够快速识别未命中模块与薄弱分支，便于进行增量补测。
  \item 局限：覆盖率反映“代码是否执行”，不直接保证业务语义正确；高覆盖不等于缺陷风险为零。
  \item 适用性：适合作为本阶段总体进度与质量基线，用于指导接口驱动用例迭代。本次测试覆盖注册登录、课程状态流转、资料上传下载、习题集维护、帖子评论等主要语句与分支。
\end{enumerate}

\subsubsection{基本路径覆盖特点}

\begin{enumerate}
  \item 优点：对登录、课程、资料、练习、社区等条件组合较多的流程有较好约束能力，能够系统性覆盖成功、失败与异常路径。
  \item 局限：路径数量随条件增加而增长较快，需要结合优先级聚焦关键业务路径，否则维护成本较高。
  \item 适用性：适合核心流程与权限相关逻辑，尤其适用于角色差异、资源存在性、状态流转等场景。
\end{enumerate}

\subsubsection{数据流分析特点}

\begin{enumerate}
  \item 优点：能够验证关键变量在跨层调用中的定义-使用一致性，擅长发现“参数传递正确但语义使用错误”的问题。
  \item 局限：需要明确变量生命周期与调用边界，分析成本高于常规接口通断测试。
  \item 适用性：适合用户身份与角色、课程标识、文件标识、密码哈希、习题集标识、题目选项与答案、检索关键字、父评论标识和匿名标记等关键变量链路验证。
\end{enumerate}

综合而言，三种方法形成互补关系：代码覆盖用于量化监控，基本路径用于保障流程完整性，数据流分析用于校验跨层变量一致性，共同支撑本阶段白盒测试结论。


\subsection{测试执行结果}

\subsubsection{执行命令与通过情况}

成员 B 执行命令：\texttt{pytest tests/ -q}

成员 B 测试结果：16 passed。

成员 C 执行命令如下：

\begin{verbatim}
pytest tests/test_exercise_module.py tests/test_post_module.py -q
\end{verbatim}

成员 C 共执行 26 个白盒测试用例，其中 23 个通过，3 个失败。3 个失败用例均为缺陷触发型用例，失败原因是当前系统缺少对应业务校验，并非测试环境或测试代码异常。

\subsubsection{覆盖率结果（本次执行）}

成员 B 覆盖率统计命令如下：

\begin{verbatim}
pytest tests/ -q --cov=controller --cov=service \
--cov=repository --cov=dao --cov-report=term-missing\end{verbatim}

成员 B 覆盖结果如下：

\begin{center}
\begin{tabular}{|p{7.4cm}|c|p{4.3cm}|}
\hline
\textbf{代码文件名} & \textbf{覆盖率} & \textbf{未覆盖行数} \\
\hline
\texttt{controller/CourseController.py} & 91\% & \texttt{18-22, 99, 114, 143} \\
\hline
\texttt{controller/LoginController.py} & 85\% & \texttt{12-16} \\
\hline
\texttt{controller/UserController.py} & 100\% & -- \\
\hline
\texttt{controller/\_\_init\_\_.py} & 100\% & -- \\
\hline
\texttt{controller/rag\_controller.py} & 100\% & -- \\
\hline
\texttt{dao/\_\_init\_\_.py} & 100\% & -- \\
\hline
\texttt{dao/document\_dao.py} & 100\% & -- \\
\hline
\texttt{repository/CourseRepository.py} & 98\% & \texttt{34} \\
\hline
\texttt{repository/UserAuthRepository.py} & 100\% & -- \\
\hline
\texttt{repository/UserRepository.py} & 92\% & \texttt{19, 29} \\
\hline
\texttt{repository/\_\_init\_\_.py} & 100\% & -- \\
\hline
\texttt{repository/document\_repository.py} & 97\% & \texttt{38} \\
\hline
\texttt{service/CourseService.py} & 100\% & -- \\
\hline
\texttt{service/LoginService.py} & 100\% & -- \\
\hline
\texttt{service/UserService.py} & 100\% & -- \\
\hline
\texttt{service/\_\_init\_\_.py} & 100\% & -- \\
\hline
\texttt{成员 B 负责代码小计} & 达到目标 & 少量非核心语句未覆盖 \\
\hline
\end{tabular}
\end{center}

成员 C 覆盖率统计命令如下：

\begin{verbatim}
pytest tests/test_exercise_module.py tests/test_post_module.py -q \
--cov=controller/ExerciseController.py --cov=controller/PostController.py \
--cov=service/ExerciseService.py --cov=service/PostService.py \
--cov=repository/ExerciseRepository.py --cov=repository/PostRepository.py \
--cov=repository/CommentRepository.py --cov-report=term-missing
\end{verbatim}

成员 C 针对练习与社区模块的覆盖率统计结果如下：

\begin{center}
\begin{tabular}{|p{7.4cm}|c|p{4.3cm}|}
\hline
\textbf{代码文件名} & \textbf{覆盖率} & \textbf{未覆盖行数} \\
\hline
\texttt{controller/ExerciseController.py} & 100\% & -- \\
\hline
\texttt{controller/PostController.py} & 100\% & -- \\
\hline
\texttt{repository/CommentRepository.py} & 100\% & -- \\
\hline
\texttt{repository/ExerciseRepository.py} & 100\% & -- \\
\hline
\texttt{repository/PostRepository.py} & 100\% & -- \\
\hline
\texttt{service/ExerciseService.py} & 100\% & -- \\
\hline
\texttt{service/PostService.py} & 100\% & -- \\
\hline
\texttt{成员 C 负责代码小计} & 100\% & -- \\
\hline
\end{tabular}
\end{center}

其中成员 C 负责模块总计 221 条语句、22 个分支，语句覆盖率和分支覆盖率均达到 100\%。该结果说明练习管理与课程社区管理中的成功路径、异常路径、权限分支、匿名/实名映射、分页过滤、评论回复等关键逻辑均已被执行。

\subsection{未覆盖代码分析}

成员 B 负责模块中仍存在少量未覆盖语句，主要并非测试遗漏或核心业务逻辑缺失，具体可分为以下三类：
\begin{itemize}
  \item \textbf{防御性分支不可触发}：当前调用链中其它层会提前进行校验，提前返回或抛出异常，导致下游用于兜底处理的部分分支不可达。例如 service 层直接抛出 \texttt{HTTPException}，该异常由 FastAPI 框架统一捕获并转换为响应，controller 层中针对返回值为 \texttt{None} 的防御分支因此无法执行。
  \item \textbf{数据库生成器语句未覆盖}：\texttt{LoginController} 与 \texttt{CourseController} 中 \texttt{get\_db()} 生成器的 \texttt{SessionLocal()} 与 \texttt{db.close()} 未被覆盖，原因是测试为保证可重复性采用 SQLite 测试会话并通过依赖覆盖替换了默认数据库依赖，从而未执行生产环境函数体。
  \item \textbf{文档存储仓储层的生命周期函数未覆盖}：仅在应用启动/关闭事件中调用，不属于常规接口请求路径。
\end{itemize}

成员 C 负责的练习与社区模块本次统计范围内未覆盖行数为 0。需要说明的是，覆盖率结果仅说明代码路径已被执行，不代表业务规则完全无缺陷；本次白盒测试正是在高覆盖执行过程中发现了 3 个输入校验与数据一致性问题。

\subsection{缺陷与问题分析}

本次白盒测试共发现 3 个明确缺陷，均来自练习与社区模块的缺陷触发型用例。

\subsubsection{系统允许上传空题目集}

\begin{enumerate}
  \item 对应测试用例：\texttt{test\_upload\_empty\_question\_list\_should\_be\_rejected}。
  \item 实际结果：接口返回 \texttt{200 OK}。
  \item 预期结果：应返回 \texttt{400} 或 \texttt{422}，拒绝空题目集。
  \item 原因分析：当前练习上传接口和习题集创建服务仅对请求进行接收和保存，没有对 \texttt{questions} 是否为空进行校验，导致系统可以创建没有题目的习题集。
\end{enumerate}

\subsubsection{系统允许对不存在的帖子发表评论}

\begin{enumerate}
  \item 对应测试用例：\texttt{test\_add\_comment\_should\_reject\_missing\_post}。
  \item 实际结果：接口返回 \texttt{200 OK}。
  \item 预期结果：应返回 \texttt{400} 或 \texttt{404}，拒绝对不存在的帖子发表评论。
  \item 原因分析：当前 \texttt{PostController.add\_comment()} 在调用评论创建逻辑前未验证 \texttt{post\_id} 是否存在，可能导致非法评论数据被接受，影响社区数据一致性。
\end{enumerate}

\subsubsection{系统允许回复不存在的父评论}

\begin{enumerate}
  \item 对应测试用例：\texttt{test\_add\_comment\_should\_reject\_missing\_parent}。
  \item 实际结果：接口返回 \texttt{200 OK}。
  \item 预期结果：应返回 \texttt{400} 或 \texttt{404}，拒绝回复不存在的父评论。
  \item 原因分析：当前 \texttt{PostController.add\_comment()} 未验证 \texttt{parent\_id} 是否对应真实评论记录，仓储层直接保存该字段，可能导致评论树出现断裂。
\end{enumerate}

\subsection{测试结论}

从白盒测试结果来看，用户与认证、课程管理、课程资料管理、练习管理和课程社区管理的主要业务路径均已得到有效验证。成员 B 负责模块的测试用例全部通过，核心业务逻辑功能正常；成员 C 负责模块的主要功能路径可正常运行，同时通过缺陷触发型用例发现了 3 个实际业务校验缺陷。

整体来看，本次白盒测试采用统一的方法与工具完成了代码覆盖、基本路径和数据流验证，主要业务分支与常见异常路径均经过验证，测试质量达到课程作业要求。对于仍未命中的少量语句，已通过代码与调用链分析确认其主要来源；对于已发现的缺陷，后续应优先补充空题目集校验、评论目标存在性校验和回复父评论存在性校验。

\section{未来扩展适配}

本部分面向练习管理与课程社区模块，基于当前实现情况推测未来可能的功能变更，并给出对应的黑盒/白盒测试调整方案。

\subsection{未来变更一：新增练习作答记录与成绩统计}

\textbf{变更描述：}
当前系统主要支持练习查看与作答反馈，未来可扩展为后端持久化作答记录、自动判题与历史成绩查询。教师端可增加按课程或习题集维度的完成率与成绩统计能力。

\textbf{黑盒测试调整：}
\begin{enumerate}
  \item 验证完整作答提交后返回分数、正确率、错题信息等结果字段。
  \item 验证重复提交规则（覆盖旧记录/保留历史/仅首次提交）符合需求定义。
  \item 验证空答案、部分答案、非法题号等异常输入的错误返回是否明确。
  \item 验证学生仅可查询本人作答记录，教师仅可查询其课程范围内统计数据。
  \item 验证未登录或无权限用户访问作答记录接口时被正确拒绝。
\end{enumerate}

\textbf{白盒测试调整：}
\begin{enumerate}
  \item 增加作答提交控制器的语句与分支覆盖（正常提交、重复提交、参数校验失败、权限拒绝）。
  \item 增加判题服务层基本路径覆盖（不同题型正确/错误/部分正确路径）。
  \item 增加成绩统计变量数据流分析（如 \texttt{score}、\texttt{correct\_count}、\texttt{wrong\_count} 的定义-使用链）。
  \item 增加作答记录仓储层插入、查询、过滤路径测试，并纳入覆盖率统计范围。
\end{enumerate}

\subsection{未来变更二：新增社区审核与删帖管理}

\textbf{变更描述：}
社区模块未来可新增管理员角色，或扩展教师管理权限，支持帖子/评论删除、内容屏蔽、举报与审核流转，以提升社区治理能力。

\textbf{黑盒测试调整：}
\begin{enumerate}
  \item 验证普通学生删除他人内容被拒绝，管理员/教师在授权范围内可执行管理操作。
  \item 验证删除、屏蔽后列表与详情页展示行为符合产品规则。
  \item 验证举报提交流程、审核状态变更与处理结果回显是否一致。
  \item 验证对已删除或已屏蔽内容继续评论/回复时，系统限制是否正确。
\end{enumerate}

\textbf{白盒测试调整：}
\begin{enumerate}
  \item 增加权限控制分支覆盖（学生无权、教师可管、管理员可管）。
  \item 增加删除逻辑基本路径覆盖（目标存在/不存在、重复删除、软删除/硬删除）。
  \item 增加举报与审核状态流的数据流分析（目标标识、处理状态、处理人字段流转）。
  \item 调整帖子/评论查询逻辑测试，验证删除标记对响应 DTO 与过滤结果的影响。
\end{enumerate}

\subsection{未来变更三：增强评论能力（多级回复与@提醒）}

\textbf{变更描述：}
在当前基础回复能力上，社区模块未来可扩展多级楼中楼、\texttt{@用户} 提醒、引用回复与评论树展示能力，并支持按时间或热度排序。

\textbf{黑盒测试调整：}
\begin{enumerate}
  \item 验证多级回复结构展示正确，父子关系和顺序符合预期。
  \item 验证 \texttt{@用户} 解析、提醒对象匹配与异常场景（不存在用户、重复提醒）处理逻辑。
  \item 验证已删除父评论、非法父节点、层级超限等边界场景返回结果。
  \item 验证评论树展开、折叠、排序切换时前后端数据一致性。
\end{enumerate}

\textbf{白盒测试调整：}
\begin{enumerate}
  \item 扩展评论创建分支覆盖（\texttt{parent\_id} 为空、合法、非法、层级超限）。
  \item 增加评论树组装基本路径覆盖（无子节点、单层、多层递归路径）。
  \item 增加提醒链路数据流分析（提醒对象解析、通知创建、事务一致性）。
  \item 调整查询断言从线性列表扩展为树形结构断言，并验证排序字段生效。
\end{enumerate}

\textbf{总体结论：}
未来扩展后，黑盒测试需重点补充业务规则、权限边界与异常输入验证；白盒测试需同步扩展控制器、服务层、仓储层的语句覆盖、分支覆盖、基本路径覆盖与关键数据流分析。建议在每次需求迭代中同步更新测试计划、测试用例、测试代码与覆盖率口径，确保扩展后系统仍具备稳定可测性。

\section{小组成员贡献与比例}

\begin{center}
\begin{tabular}{|p{1.8cm}|p{8.8cm}|c|}
\hline
\textbf{成员} & \textbf{主要贡献} & \textbf{贡献比例} \\
\hline
桂舒芸（成员 A） & 负责黑盒测试整体工作，包括测试范围分析、黑盒测试方法选择、测试用例设计、自动化接口测试代码编写、人工页面验证、测试执行结果整理和黑盒测试报告撰写。 & 33.3\% \\
\hline
孙昕玥（成员 B） & 负责用户与认证、课程管理、课程资料管理相关模块的白盒测试，包括代码结构分析、覆盖率测试、基本路径与数据流分析、测试代码编写、执行结果统计和报告整理。 & 33.3\% \\
\hline
石欣睿（成员 C） & 负责练习管理与课程社区管理相关模块的白盒测试，包括代码覆盖、基本路径、数据流分析、缺陷触发型用例设计、覆盖率结果整理，并补充未来扩展适配方案。 & 33.3\% \\
\hline
\end{tabular}
\end{center}

\end{document}
